\documentclass[11pt]{article}
\usepackage{amsmath,amssymb,amsthm}
\newtheorem{theorem}{Theorem}
\title{Exact Parity Correlation of Bounded Admissible Atoms}
\author{}
\date{November 23, 2026}
\begin{document}
\maketitle

\paragraph{Status and scope.}
This is a complete written proof, with an exact finite arithmetic check,
but not a Lean formalization. It concerns an individually bounded atom.
It is not a lower bound for arbitrary sums of atoms, because their
individual magnitudes can be unbounded while their sum remains bounded.
In particular it establishes neither of the requested global asymptotic
alternatives. No model definition is changed.

\begin{theorem}
Let $n\geq1$, let $b_0>0$, and let $b_i\geq0$. Put
\[
 B(x)=b_0+\sum_{i=1}^n b_ix_i,\qquad
 A(x)=a_0+\sum_{i=1}^n a_ix_i.
\]
If $|A(x)/B(x)|\leq1$ at every Boolean vertex, then
\[
 \left|2^{-n}\sum_{x\in\{0,1\}^n}
              (-1)^{|x|}\frac{A(x)}{B(x)}\right|\leq 2^{1-n}.
 \tag{1}
\]
The same bound holds when $B$ is positive everywhere on the cube and
all its slopes are nonpositive. For strictly positive slopes, the
constant in (1) is a supremum and cannot be improved.
\end{theorem}
\begin{proof}
Write
\[
 I_0=\int_0^\infty e^{-b_0t}\prod_j(1-e^{-b_jt})\,dt,
 \qquad
 I_i=\int_0^\infty e^{-(b_0+b_i)t}
                    \prod_{j\ne i}(1-e^{-b_jt})\,dt.
\]
All these integrals are nonnegative and finite. Expanding the finite
products and using $\int_0^\infty e^{-ct}\,dt=1/c$ for $c>0$ gives
\[
 \sum_x(-1)^{|x|}\frac{A(x)}{B(x)}
           =a_0I_0-\sum_i a_i I_i.
 \tag{2}
\]
Integration of the derivative of
$e^{-b_0t}\prod_j(1-e^{-b_jt})$ gives
\[
                 \sum_i b_iI_i=b_0I_0.
 \tag{3}
\]
Indeed its values at both endpoints are zero, using $n\geq1$ at
$t=0$. This also covers zero slopes. Furthermore, for every $i$,
\[
 I_0+I_i
 =\int_0^\infty e^{-b_0t}\prod_{j\ne i}(1-e^{-b_jt})\,dt
 \leq\frac1{b_0}.
 \tag{4}
\]

Set $d_i=a_i+b_i$ and $c=a_0+b_0$. The lower bound $A+B\geq0$
at every vertex implies
\[
             \sum_i\max\{-d_i,0\}\leq c.
 \tag{5}
\]
This follows by choosing exactly the negative $d_i$ coordinates to be
one. The bounds at the zero vertex give $0\leq c\leq2b_0$.
Combining (2) and (3), and putting $M=\max_i I_i$, now yields
\[
 \begin{aligned}
 a_0I_0-\sum_i a_iI_i
 &=cI_0-\sum_i d_iI_i\\
 &\leq cI_0+M\sum_i\max\{-d_i,0\}\\
 &\leq c(I_0+M)\leq 2.
 \end{aligned}
\]
Equation (4) proves the last step. Replacing $A$ by $-A$ proves the
opposite inequality. Division by $2^n$ gives (1).

If every slope is nonpositive, complement every input coordinate.
The new denominator has positive constant coefficient $B(\mathbf1)$
and nonnegative slopes. Complementation changes parity by the fixed
factor $(-1)^n$, so its absolute correlation is unchanged.

For sharpness take $b_0=1$, $b_i=T>0$ for every $i$, and
\[
 B_T(x)=1+T|x|,\qquad
 A_T(x)=1-T|x|-2x_1.
\]
Then $A_T+B_T=2(1-x_1)\geq0$ and
$B_T-A_T=2T|x|+2x_1\geq0$, so the quotient is bounded by one.
As $T\to\infty$ it tends to $+1$ at the zero vertex and $-1$ at
every other vertex. Its unnormalized parity correlation therefore
tends to $2$, proving sharpness. Each denominator has strictly
positive slopes, so this family consists of original admissible
atoms via the existing positive-affine-ratio construction.
\end{proof}

\paragraph{Consequence for the proposed approximation route.}
A distribution-uniform weak-learning claim with correlation
$\Omega(n2^{-n/2})$ for individually bounded admissible heads is false:
parity under the uniform distribution violates it by (1). The same
holds for any convex combination of such heads and for its negative.
An unrestricted affine numerator or an unrestricted outer coefficient
does not contradict the theorem; it invalidates the individual
boundedness hypothesis unless the resulting head is normalized again.
The theorem does not restrict a sum whose heads have large cancelling
magnitudes. A global improvement, if it exists, may use precisely that
behavior.

\end{document}
